% Les trois modes de transfert thermique (L'essentiel, chapitre 20)
\begin{tikzpicture}[x=1cm,y=1cm]
% conduction : barre chauffée à une extrémité
\shade[left color=otRed!70,right color=otNavy!35] (0,1.0) rectangle (3.2,1.45);
\draw[otInk!70,line width=0.5pt] (0,1.0) rectangle (3.2,1.45);
\foreach \xx in {0.4,1.1,1.8}{\draw[otOrange!90!black,line width=0.8pt,-{Stealth[length=1.5mm]}] (\xx,1.7) -- (\xx+0.7,1.7);}
\fill[otOrange!80] (0.05,0.25) .. controls (0.25,0.7) and (0.35,0.55) .. (0.3,0.95) .. controls (0.55,0.6) and (0.55,0.4) .. (0.55,0.25) -- cycle;
\node[pcLegendeG,anchor=north] at (1.6,-0.05) {conduction};
\node[pcLegende,anchor=north,align=center] at (1.6,-0.42) {de proche en proche,\\ sans déplacement de matière};
% convection : casserole
\begin{scope}[shift={(4.3,0)}]
\fill[pcEau] (0.1,0.45) rectangle (2.7,2.0);
\draw[otInk!70,line width=0.8pt] (0.1,2.2) -- (0.1,0.45) -- (2.7,0.45) -- (2.7,2.2);
\draw[otRed!75!black,line width=0.8pt,-{Stealth[length=1.6mm]}] (1.4,0.65) .. controls (1.4,1.5) .. (1.4,1.75) .. controls (1.4,1.9) and (0.6,1.9) .. (0.5,1.6);
\draw[otNavy,line width=0.8pt,-{Stealth[length=1.6mm]}] (0.5,1.55) .. controls (0.45,0.8) and (0.8,0.65) .. (1.3,0.65);
\draw[otRed!75!black,line width=0.8pt,-{Stealth[length=1.6mm]}] (1.4,1.75) .. controls (1.4,1.9) and (2.2,1.9) .. (2.3,1.6);
\draw[otNavy,line width=0.8pt,-{Stealth[length=1.6mm]}] (2.3,1.55) .. controls (2.35,0.8) and (2.0,0.65) .. (1.5,0.65);
\fill[otOrange!80] (0.9,0.0) rectangle (1.9,0.3);
\node[pcLegendeG,anchor=north] at (1.4,-0.05) {convection};
\node[pcLegende,anchor=north,align=center] at (1.4,-0.42) {déplacement\\ d'un fluide};
\end{scope}
% rayonnement : soleil et corps éclairé
\begin{scope}[shift={(8.2,0)}]
\fill[otOrange!80] (0.45,1.3) circle (0.4);
\foreach \a in {0,45,...,315}{\draw[otOrange!80,line width=0.6pt] ($(0.45,1.3)+(\a:0.47)$) -- ($(0.45,1.3)+(\a:0.62)$);}
\foreach \yy in {0.9,1.3,1.7}{\draw[otOrange!90!black,line width=0.8pt,-{Stealth[length=1.5mm]},decorate,
  decoration={snake,amplitude=0.6mm,segment length=2.6mm,post length=2mm}] (1.2,\yy) -- (2.45,\yy);}
\fill[otMuted!55] (2.55,0.6) rectangle (2.95,2.0);
\node[pcLegendeG,anchor=north] at (1.6,-0.05) {rayonnement};
\node[pcLegende,anchor=north,align=center] at (1.6,-0.42) {ondes électromagnétiques,\\ même dans le vide};
\end{scope}
\end{tikzpicture}
\par\smallskip\legende{Un transfert thermique peut s'effectuer par conduction, par convection ou par rayonnement ; les trois modes coexistent souvent.}
